\section{Original source, Burnol's division functions and local factors}
Let $D=-x\partial_x$, $\mathcal H=L^2(\mathbb R_{>0},dx)$ and
\[
 g(s)=2\xi(s)=s(s-1)L_\infty(s)\zeta(s),\quad
 L_\infty(s)=\pi^{-s/2}\Gamma(s/2),\quad
 h(s)=\prod_{\rho\in Z}(s-\rho)^{m_\rho},\quad d=\deg h.\tag{LR1}
\]
Here $Z$ is a fixed finite set of nontrivial zeros and each $m_\rho$ is
its full multiplicity. The original source $a_h$ has Mellin transform
$g/h$. Write $\mathcal M f(s)=\int_0^\infty f(x)x^{s-1}dx$,
$j_hF=(F^{(r)}(\rho)/r!)_{\rho,r<m_\rho}$, and
$\mathsf T_h=j_h(g/h)\cdot$, multiplication in the full jet algebra.
It is invertible. The physical source and its jet are exactly
\[
 \begin{gathered}
 \mathcal A_NP=P(D)a_h,\quad j_h\mathcal M(\mathcal A_NP)
 =\mathsf T_hj_hP,\\
 \|\mathcal A_NP\|^2=\int_{\mathbb R}|P(1/2+iy)|^2w_h(y)dy,
 \quad w_h(y)=\frac{|(g/h)(1/2+iy)|^2}{2\pi}.
 \end{gathered}\tag{LR2}
\]
Every formula below uses this measure including its $2\pi$ factor.

\begin{proposition}[Burnol-to-original-source map, with full divided jets]
Let $P$ have degree below $d$. Its exact partial fraction decomposition is
\[
 \frac{P(s)}{h(s)}=\sum_{\rho\in Z}\sum_{l=1}^{m_\rho}
 \frac{c_{\rho l}(P)}{(s-\rho)^l},\quad
 c_{\rho l}(P)=\frac1{(m_\rho-l)!}
 \left[\partial_s^{m_\rho-l}\frac{P(s)}{h(s)/(s-\rho)^{m_\rho}}
 \right]_{s=\rho}.\tag{LR3}
\]
There is a unique $b_P\in L^2(0,\infty;dx)$ with right Mellin transform
$\mathcal M_Rb_P(s)=\zeta(s)P(s)/h(s)$. It belongs to Burnol's
$L_1$ (both it and its cosine transform are constant on $(0,1)$).
Define $If(x)=x^{-1}f(1/x)$ and the bounded Mellin multiplier
$T_m=\mathcal M^{-1}m(s)\mathcal M$, $m(s)=s(s-1)L_\infty(s)$,
on the line $\Re s=1/2$. Then
\[
 \mathcal A_{d-1}P=T_m I b_P,\qquad
 \|\mathcal A_{d-1}P\|^2=
 \int_{\mathbb R}|m(1/2+iy)|^2
 |\mathcal M_Rb_P(1/2+iy)|^2\frac{dy}{2\pi}.\tag{LR4}
\]
Thus the original canonical section is obtained by inserting
$P=r_h(u)$, $j_hr_h(u)=\mathsf T_h^{-1}u$ in (LR3)--(LR4).
\end{proposition}
\begin{proof}
Principal parts give (LR3); after subtracting them the rational function
has no poles and vanishes at infinity, so it is zero. Burnol
\cite[\S4, Proposition \texttt{zetaprop} and its proof]{Burnol}
proves that each $\zeta(s)/(s-\rho)^l$, $l\le m_\rho$, is the
right Mellin transform of a vector of $L_1$. Its proof uses the exact
Hardy-space characterization of $L_1$ and
$\chi(s)\zeta(1-s)/(1-s-\rho)^l
=(-1)^l\zeta(s)/(s-(1-\rho))^l$; it does not assume RH.
The finite linear combination (LR3) therefore defines $b_P$.
Mellin Plancherel proves uniqueness.
Substitution $t=1/x$ shows both $\|If\|=\|f\|$ and
$\mathcal M(If)(s)=\mathcal M_Rf(s)$. The function $m$ is continuous
on this line and bounded: the Gamma asymptotic
\cite[5.11.9]{DLMF5} gives polynomial growth times
$e^{-\pi|y|/4}$ at either end. Consequently $T_m$ is bounded.
Its product with $\zeta P/h$ is exactly $gP/h$, proving (LR4) by
Plancherel. The multiplier has no zero on the line, so the map is
injective. Its inverse on this finite-dimensional image is the
restriction of multiplication by $1/m$; no unweighted isometry is asserted.
\end{proof}

For the full original polynomial space, division gives a unique
$P=r+hQ$, $\deg r<d$. Consequently
\[
 \mathcal A_NP=T_m I b_r+Q(D)f_0,\qquad
 \mathcal M f_0=g,\qquad h(D)a_h=f_0.\tag{LR5}
\]
This explicitly retains every relation vector and its cross terms in
the norm. It does not assert that $\zeta Q$ is an $L^2$ Mellin transform.
The full quotient minimum below includes the complete second summand.

\begin{proposition}[Exact semilocal moment receiver]
For a finite set $S$ of places containing $\infty$, put
\[
 a_S(y)=L_\infty(s)\prod_{p\in S\setminus\{\infty\}}
 (1-p^{-s})^{-1},\quad s=1/2+iy,\quad d\mu_S(y)=|a_S(y)|^2dy,
\]
\[
 q_{S,h}(y)=\frac{g(s)}{\sqrt{2\pi}h(s)a_S(y)}
 =\frac{s(s-1)\zeta(s)}{\sqrt{2\pi}h(s)}
 \prod_{p\in S\setminus\{\infty\}}(1-p^{-s}).\tag{LR6}
\]
Multiplication $U_{S,h}F=q_{S,h}F$ is unitary from
$L^2(w_hdy)$ to $L^2(d\mu_S)$. It intertwines multiplication by $y$
on their full maximal domains. The native polynomial space maps exactly
to $q_{S,h}\mathcal P_N$, and its norm, restrictions, full relation
minima and attained observation quotients are preserved under this map.
\end{proposition}
\begin{proof}
The density identity $|q_{S,h}|^2|a_S|^2=w_h$ proves isometry.
The finite Euler factors have no zero on the line. The quotient $g/h$
is entire and its zeros there are discrete, so $q_{S,h}\ne0$ almost
everywhere. For any $V\in L^2(d\mu_S)$, the function $V/q_{S,h}$
has squared $w_h$ norm equal to $\|V\|^2$, proving surjectivity.
The same identity with $|yF|^2$ proves the domain and intertwining
assertion. Applying this identity to each vector in every original
affine fibre preserves its infimum and its minimizing vector.
It also preserves each cross inner product by polarization.
\end{proof}

Connes--Consani--Moscovici \cite[\S3.2, Proposition 3.1]{CCM}
prove determinacy for precisely this local-factor density (their
choice of sign on the line leaves the modulus unchanged).
Their Gamma estimate bounds its even moments by $M^n(2n)!$.
No division by total mass has occurred in (LR6). To make its Jacobi
form explicit, let $V_S:L^2(\mu_S)\longrightarrow\ell^2$ be expansion
in its orthonormal polynomials, and put
$J_S=V_SM_yV_S^{-1}$ and $\xi_S=\sqrt{\mu_S(\mathbb R)}e_0$.
Polynomial completeness here follows, for example, directly from
the exponential tail: if $F$ is orthogonal to every polynomial,
the Fourier transform of $F\mu_S$ is analytic in a strip by
Cauchy--Schwarz, has every derivative zero at zero, and hence vanishes;
Fourier uniqueness gives $F=0$. Thus $V_S$ is unitary. For the original
polynomials $P_i$ one has the exact full Gram
\[
 G_{ij}=\left\langle q_{S,h}(J_S)P_i(1/2+iJ_S)\xi_S,
 q_{S,h}(J_S)P_j(1/2+iJ_S)\xi_S\right\rangle.\tag{LR7}
\]
The functional-calculus vectors are in their stated domains because
(LR2) is finite. The root polynomial, zeta factor and finite Euler
factors in (LR6) remain inside this expression. In particular no
finite Euler-product limit on the critical line is used.

\section{Evaluating the original interpolation minimum by its kernel}
All moments of $w_h$ are finite by the Gamma decay in (LR1), and its
support is infinite. Let $p_n$ be its real monic orthogonal polynomials,
$h_n=\int p_n(y)^2w_h(y)dy>0$. These are auxiliary coordinates in the
unchanged source space, not a replacement of its measure. They can be
computed without an unknown scalar: with
$c_j=\int y^jw_h(y)dy$ and
$\Delta_{n-1}=\det(c_{i+j})_{0\le i,j<n}$, $\Delta_{-1}=1$,
\[
 p_n(z)=\frac1{\Delta_{n-1}}
 \det\begin{pmatrix}
 c_0&c_1&\cdots&c_n\\
 c_1&c_2&\cdots&c_{n+1}\\
 \vdots&\vdots&&\vdots\\
 c_{n-1}&c_n&\cdots&c_{2n-1}\\
 1&z&\cdots&z^n
 \end{pmatrix},\qquad h_n=\Delta_n/\Delta_{n-1}.\tag{LR8}
\]
Expansion of the determinant shows it is monic and orthogonal to the
first $n$ monomials; the norm identity follows by block elimination.
This is the moment construction used in \cite[\S2.1]{CCM}, with all
mass factors retained.

\begin{proposition}[Full-jet Christoffel--Darboux minimum]
For $N\ge d-1$ set
\[
 K_N(z,w)=\sum_{n=0}^N\frac{p_n(z)\overline{p_n(w)}}{h_n}
 =\frac{p_{N+1}(z)\overline{p_N(w)}-
 p_N(z)\overline{p_{N+1}(w)}}{h_N(z-\bar w)}.\tag{LR9}
\]
At $z=\bar w$ use the removable value and its derivatives.
For each $n$, retain the complete physical-jet vector
\[
 t_n=j_h\left[(g/h)(s)\frac{p_n((s-1/2)/i)}{\sqrt{h_n}}\right],
 \qquad M_N=\sum_{n=0}^N t_nt_n^*.\tag{LR10}
\]
Then $M_N>0$, and the original attained metric is
\[
 H_N=M_N^{-1},\qquad
 \min_{\mathsf T_hj_hP=u}\|\mathcal A_NP\|^2=u^*H_Nu.
 \tag{LR11}
\]
Thus every matrix entry in $M_N$ is obtained by taking the full divided
derivatives in (LR9), including the $g/h$ factors and $i^{-r}$ factors
from $s=1/2+iy$. The exact next-cutoff update is
\[
 H_{N+1}=H_N-\frac{H_Nt_{N+1}t_{N+1}^*H_N}
 {1+t_{N+1}^*H_Nt_{N+1}},\quad
 \log\frac{\det H_{N+1}}{\det H_N}
 =-\log(1+t_{N+1}^*H_Nt_{N+1}).\tag{LR12}
\]
\end{proposition}
\begin{proof}
Symmetry of multiplication by $y$ gives
$yp_n=p_{n+1}+\alpha_np_n+(h_n/h_{n-1})p_{n-1}$; coefficients
of all lower $p_j$ vanish by orthogonality. Multiply this recurrence
at $z$ by $\overline{p_n(w)}/h_n$, subtract its conjugate at $w$,
and sum. All internal terms cancel, leaving (LR9). This proves the
classical Christoffel--Darboux formula \cite[18.2.12]{DLMF18}
also at complex points by polynomial continuation.

Write $P(s)=\sum_{n=0}^N c_np_n((s-1/2)/i)/\sqrt{h_n}$.
Then its exact squared source norm is $c^*c$ and its physical jet is
$T_Nc$, $T_N=(t_0,\ldots,t_N)$. Hermite interpolation and the unit
$\mathsf T_h$ imply that $T_N$ is onto for $N\ge d-1$.
The coefficient vector of least norm on $T_Nc=u$ is
$c_*=T_N^*(T_NT_N^*)^{-1}u$: it has the required image and is
orthogonal to $\ker T_N$. Every other vector is $c_*+v$ with
$v\in\ker T_N$, so its squared norm is $\|c_*\|^2+\|v\|^2$.
This proves (LR11). In polynomial coordinates the same kernel is
exactly $h\mathcal P_{N-d}$, including all terms in (LR5).
Formula (LR10) gives $M_{N+1}=M_N+tt^*$; multiplication verifies
the inverse formula (LR12), and the rank-one determinant identity
gives its second formula. No dimension factor was discarded.
\end{proof}

For the tensor-sum source, use its unchanged measure $w_h^{*k}dy$,
centre $c=k/2$, and full sum-root polynomial $\chi_k$ in the same
orthogonal-polynomial construction. In its original polynomial-class
coordinates set $t_n=j_{\chi_k}[p_n((S-c)/i)/\sqrt{h_n}]$; if the
physical tensor coordinates include a further jet unit, multiply by
that exact unit. Hermite interpolation again proves (LR10)--(LR12).
This first gives the attained class metric $G_N=H_N$ in the specified
coordinates. For the original observation $\Lambda$ on this class
space, the second minimum is
$Q_N=(\Lambda H_N^{-1}\Lambda^*)^{-1}$ on its image. Thus the
construction gives entries for the original class metric, rather than
confusing it with the larger coefficient Gram before the first minimum.
Insert it in $\Omega=\Lambda^*Q_N\Lambda$, $L=H_N-\Omega$ and retain
the actual class vectors $a_i,v$ to compute
$K_i=a_i^*Lv$, $B_i=a_i^*\Omega v$. Their HG7 signed products then
keep both cross terms. This is a proved receiving calculation;
their growing-degree asymptotics and individual signs have not been
evaluated by replacing $q_{S,h}\mathcal P_N$ by $\mathcal P_N$.

\section{Exact source correspondence with the classical zeta heat flow}
The original theta source is
\[
 \phi_*(x)=(4\pi^2x^4-6\pi x^2)e^{-\pi x^2},\qquad
 f_0(x)=2\sum_{n\ge1}\phi_*(nx),\qquad \mathcal M f_0=g.
\]
For Polymath's $\Phi$ and $H_t$ \cite[equations (1)--(4)]{Polymath},
direct substitution, with every numerical factor, gives
\[
 \Phi(u)=\tfrac14e^u f_0(e^{2u}),\quad
 f_t(x)=e^{t(\log x)^2/4}f_0(x),\quad
 \mathcal M f_t(1/2+iy)=16H_t(2y).\tag{LR13}
\]
Indeed the two coefficients become $2\pi^2n^4e^{9u}$ and
$-3\pi n^2e^{5u}$. On substituting $x=e^{2u}$ in the Mellin
integral its Jacobian and $x^{1/2}$ give $8\Phi(u)du$.
The evenness of $\Phi$ then gives the remaining factor two.
The integrals and all derivatives converge by the two-sided
superexponential decay supplied by theta inversion. Consequently
\[
 \partial_t\mathcal M f_t(s)=\tfrac14\partial_s^2\mathcal M f_t(s),
 \quad \partial_tH_t=-\partial_z^2H_t.\tag{LR14}
\]
If $B_t$ denotes multiplication by $e^{t(\log x)^2/4}$ on its domain,
direct differentiation gives the exact conductor transformation
\[
 h(D)B_t=B_t\,h(D-\tfrac t2\log x),\qquad
 [D,\log x]=-I.\tag{LR15}
\]
The powers in the right side are operator compositions in their
written order; commuting the factors would change the formula.
For a simple zero $\rho(t)$ of $g_t=\mathcal M f_t$, the implicit
function theorem gives its actual velocity
\[
 \rho'(t)=-\frac{g_t''(\rho(t))}{4g_t'(\rho(t))}.\tag{LR16}
\]
Multiple roots continue to require the full primary algebra; (LR16)
is asserted only on the simple-root domain. Connes's multiplicative
Gaussian in \cite{ConnesHeat} instead acts by
$e^{t(s-1/2)^2}$ on the Mellin transform. Its physical convolution
and the Gram transport are proved in GM below. Equations
(LR13)--(LR15) give the exact relation of the two operations to the
same original source and conductor, rather than identifying their
different parameter evolutions.

\section{Marked Gaussian formula with its arithmetic terms retained}
For $t>0$ and $a\in\mathbb R$, let
$k_t(v)=e^{-v^2/(4t)}/\sqrt{4\pi t}$ and
$F_{t,a}(x)=k_t(\log x-a)$, $f_{t,a}=x^{-1/2}F_{t,a}$.
Connes's explicit formula \cite[\S2, equations (4)--(12)]{ConnesHeat}
applies to this rapidly decreasing test function. Its exact marked form is
\[
 \begin{split}
 \sum_\rho m_\rho e^{t(\rho-1/2)^2+a(\rho-1/2)}
 ={}&2e^{t/4}\cosh(a/2)-W_{\mathbb R}(F_{t,a})\\
 &-\sum_{n\ge2}\frac{\Lambda(n)}{\sqrt n}
 [k_t(\log n-a)+k_t(\log n+a)],
 \end{split}\tag{LR17}
\]
where the sum on the left ranges over all distinct nontrivial zeros,
and $\Lambda$ is the von Mangoldt function. The full archimedean term is
\[
 \begin{split}
 W_{\mathbb R}(F)={}&(\log4\pi+\gamma_E)F(1)\\
 &+\int_0^\infty
 [F(e^u)+F(e^{-u})-2e^{-u/2}F(1)]
 \frac{e^{u/2}}{e^u-e^{-u}}du .
 \end{split}\tag{LR18}
\]
To prove (LR17), substitute $x=e^v$ and complete the square to obtain
$\mathcal M f_{t,a}(s)=e^{t(s-1/2)^2+a(s-1/2)}$.
The two pole terms are its values at one and zero. The prime-power terms
in Connes's formula are exactly the displayed sum over $n$, and
$x=e^u$ gives (LR18). This proof does not require RH or a
self-adjoint zero operator. The zero sum is absolutely and locally
uniformly convergent for $t>0$: $0<\Re\rho<1$, Gaussian damping in
$\Im\rho$, and the classical $O(T\log(T+2))$ zero count suffice.
The prime sum has the same property by Gaussian decay in $\log n$.
All fixed-order $a$ derivatives therefore pass through both sums and
the integral, by domination with a Gaussian times a polynomial.
For example
\[
 \partial_a^r k_t(u-a)=
 (2\sqrt t)^{-r}\mathcal H_r((u-a)/(2\sqrt t))k_t(u-a),\tag{LR19}
\]
where $\mathcal H_r$ is the physicists' Hermite polynomial, as follows
by differentiating the Gaussian. The derivative of $k_t(u+a)$ gains
$(-1)^r$ relative to the same expression. Formula (LR19) keeps the
precise growing-order factors. Fixed-order flatness of the prime term
as $t\downarrow0$ does not estimate it after order $r$ grows with
the programme degree. For a finite full jet algebra the corresponding
marked operator is exactly
$j_h\exp(t(s-1/2)^2+a(s-1/2))\cdot$; its nilpotent terms must be
kept when applying (LR17) to marked sources.
